\section{Giới hạn của hàm số}

\subsection{Đề 1}
\Opensolutionfile{ans}[ans/ans-1-C3B2]
\caulc
\begin{ex}%[1D3N2-2]%[Dự án tex tài liệu Dương Hùng-Ngô Tất Thành]
	Giá trị của giới hạn $\lim\limits_{x\to 2}(3x^2+7x+11)$ là
	\choice
	{\True $37$}
	{$38$}
	{$39$}
	{$40$}
	\loigiai{
		Ta có $\lim\limits_{x\to 2}(3x^2+7x+11)=3 \cdot 2^2+7 \cdot 2+11=37$.
	}
\end{ex}

\begin{ex}%[1D3N2-2]%[Dự án tex tài liệu Dương Hùng-Ngô Tất Thành]
	Giá trị của giới hạn $\lim\limits_{x\to\sqrt 3}|x^2-4|$ là
	\choice
	{$0$}
	{\True $1$}
	{$2$}
	{$3$}
	\loigiai{
		Ta có $\lim\limits_{x\to\sqrt 3}|x^2-4|=|(\sqrt 3)^2-4|=1$.
	}
\end{ex}

\begin{ex}%[1D3N2-2]%[Dự án tex tài liệu Dương Hùng-Ngô Tất Thành]
	Giá trị của giới hạn $\lim\limits_{x\to 0}{x^2}\sin\dfrac{1}{2}$ là
	\choice
	{$\sin\dfrac{1}{2}$}
	{$+\infty$}
	{$-\infty$}
	{\True $0$}
	\loigiai{
		Ta có $\lim\limits_{x\to 0}{x^2}\sin\dfrac{1}{2}=0 \cdot \sin\dfrac{1}{2}=0$.
	}
\end{ex}

\begin{ex}%[1D3N2-2]%[Dự án tex tài liệu Dương Hùng-Ngô Tất Thành]
	Giá trị của giới hạn $\lim\limits_{x\to -1}\dfrac{x^2-3}{x^3+2}$ là
	\choice
	{$1$}
	{\True $-2$}
	{$2$}
	{$-\dfrac{3}{2}$}
	\loigiai{
		Ta có  $\lim\limits_{x\to-1}\dfrac{x^2-3}{x^3+2}= \dfrac{(-1)^2-3}{(-1)^3+2}=-2$.
	}
\end{ex}

\begin{ex}%[1D3N2-2]%[Dự án tex tài liệu Dương Hùng-Ngô Tất Thành]
	Giá trị của giới hạn $\lim\limits_{x\to 1}\dfrac{x-x^3}{(2x-1)(x^4-3)}$ là
	\choice
	{$1$}
	{$-2$}
	{\True $0$}
	{$-\dfrac{3}{2}$}
	\loigiai{
		Ta có  $\lim\limits_{x\to 1}\dfrac{x-x^3}{(2x-1)(x^4-3)}=
		\dfrac{1-1^3}{(2 \cdot 1-1)(1^4-3)}=0$.
	}
\end{ex}

\begin{ex}%[1D3N2-2]%[Dự án tex tài liệu Dương Hùng-Ngô Tất Thành]
	Giá trị của giới hạn $\lim\limits_{x\to -1}\dfrac{|x-1|}{x^4+x-3}$ là
	\choice
	{$-\dfrac{3}{2}$}
	{$\dfrac{2}{3}$}
	{$\dfrac{3}{2}$}
	{\True $-\dfrac{2}{3}$}
	\loigiai{
		Ta có $\lim\limits_{x\to-1}\dfrac{|x-1|}{x^4+x-3}= \dfrac{|-1-1|}{1-1-3}=-\dfrac{2}{3}$.
	}
\end{ex}

\begin{ex}%[1D3N2-2]%[Dự án tex tài liệu Dương Hùng-Ngô Tất Thành]
	Giá trị của giới hạn $\lim\limits_{x\to 3}\sqrt{\dfrac{9x^2-x}{(2x-1)(x^4-3)}}$ là
	\choice
	{$\dfrac{1}{5}$}
	{$\sqrt 5$}
	{\True $\dfrac{1}{\sqrt 5}$}
	{$5$}
	\loigiai{
		Ta có $\lim\limits_{x\to 3}\sqrt{\dfrac{9x^2-x}{(2x-1)(x^4-3)}}
		=\sqrt{\dfrac{9 \cdot 3^2-3}{(2 \cdot 3-1)(3^4-3)}}
		=\dfrac{1}{\sqrt 5}$.
	}
\end{ex}

\begin{ex}%[1D3N2-2]%[Dự án tex tài liệu Dương Hùng-Ngô Tất Thành]
	Giá trị của giới hạn $\lim\limits_{x\to 2}\sqrt[3]{\dfrac{x^2-x+1}{x^2+2x}}$ là
	\choice
	{$\dfrac{1}{4}$}
	{\True $\dfrac{1}{2}$}
	{$\dfrac{1}{3}$}
	{$\dfrac{1}{5}$}
	\loigiai{
		Ta có $\lim\limits_{x\to 2}\sqrt[3]{\dfrac{x^2-x+1}{x^2+2x}}= \sqrt{\dfrac{2^2-2+1}{2^2+2 \cdot 2}}=\dfrac{1}{2}$.
	}
\end{ex}

\begin{ex}%[1D3H2-2]%[Dự án tex tài liệu Dương Hùng-Ngô Tất Thành]
	Giá trị của giới hạn $\lim\limits_{x\to 2}\dfrac{\sqrt[3]{3x^2-4}-\sqrt{3x-2}}{x+1}$ là
	\choice
	{$-\dfrac{3}{2}$}
	{$-\dfrac{2}{3}$}
	{\True $0$}
	{$+\infty$}
	\loigiai{
		Ta có $\lim\limits_{x\to 2}\dfrac{\sqrt[3]{3x^2-4}-\sqrt{3x-2}}{x+1}= \dfrac{\sqrt[3]{12-4}-\sqrt{6-2}}{3}=\dfrac{0}{3}=0$.
	}
\end{ex}

\begin{ex}%[1D3H2-4]%[Dự án tex tài liệu Dương Hùng-Ngô Tất Thành]
	Giá trị của giới hạn $\lim\limits_{x\to+\infty}\left(\sqrt{x^2+1}+x\right)$ là
	\choice
	{$0$}
	{\True $+\infty$}
	{$\sqrt 2-1$}
	{$-\infty$}
	\loigiai{
		Ta có $\lim\limits_{x\to+\infty}\left(\sqrt{x^2+1}+x\right)= \lim\limits_{x\to+\infty}x\left(\sqrt{1+\dfrac{1}{x^2}}+1\right)= +\infty$.\\
		vì $\heva{
			&\lim\limits_{x\to+\infty}x=+\infty\\
			&\lim\limits_{x\to{+\infty}}\sqrt{1+\dfrac{1}{x^2}}+1=2 > 0.}$
	}
\end{ex}

\begin{ex}%[1D3H2-4]%[Dự án tex tài liệu Dương Hùng-Ngô Tất Thành]
	Giá trị của giới hạn $\lim\limits_{x\to+\infty}\left(\sqrt[3]{3x^3-1}+\sqrt{x^2+2}\right)$ là
	\choice
	{$\sqrt[3]{3}+1$}
	{\True $+\infty$}
	{$\sqrt[3]{3}-1$}
	{$-\infty$}
	\loigiai{
		Ta có 
		$\lim\limits_{x\to+\infty}\left(\sqrt[3]{3x^3-1}+\sqrt{x^2+2}\right)= \lim\limits_{x\to+\infty}x\left(\sqrt[3]{3-\dfrac{1}{x^3}}+\sqrt{1+\dfrac{2}{x^2}}\right)=+\infty$.
	}
\end{ex}

\begin{ex}%[1D3H2-7]%[Dự án tex tài liệu Dương Hùng-Ngô Tất Thành]
	Kết quả của giới hạn $\lim\limits_{x\to{2^+}}\dfrac{\sqrt{x+2}}{\sqrt{x-2}}$ là
	\choice
	{$-\infty$}
	{\True $+\infty$}
	{$-\dfrac{15}{2}$}
	{Không xác định}
	\loigiai{
		Ta có $\heva{
			&\lim\limits_{x\to{2^+}}\sqrt{x+2}=2>0\\
			&\lim\limits_{x\to{2^+}}\sqrt{x-2}=0}$
		và $\sqrt{x-2}>0$, $\forall x>2$.\\
		Suy ra $\lim\limits_{x\to{2^+}}\dfrac{\sqrt{x+2}}{\sqrt{x-2}}=+\infty.$
	}
\end{ex}
\Closesolutionfile{ans}
\indapan{6}{ans/ans-1-C3B2}

\cauds
\Opensolutionfile{ans}[ans/ans-1-C3B2-DS]
\begin{ex}%[1D3H2-3]%[Dự án tex tài liệu Dương Hùng-Ngô Tất Thành]
	Xét tính đúng sai của các mệnh đề sau
	\choiceTF
	{\True $\lim\limits_{x\to-2}(x^2-x+3)=9$}
	{$\lim\limits_{x\to 6}\sqrt{\dfrac{1}{x+3}}=3$}
	{\True $\lim\limits_{x\to 2}\dfrac{x^2-3x+2}{x-2}=1$}
	{$\lim\limits_{x\to-1}\dfrac{2x^2+3x+1}{x^2-1}=\dfrac{1}{3}$}
	\loigiai{
		\begin{itemchoice}
			\itemch Ta có $\lim\limits_{x\to-2}(x^2-x+3)=(-2)^2-(-2)+3=9$.
			\itemch Ta có $\lim\limits_{x\to 6}\sqrt{\dfrac{1}{x+3}}=\sqrt{\dfrac{1}{6+3}}=\dfrac{1}{3}$.
			\itemch Ta có $\lim\limits_{x\to 2}\dfrac{x^2-3x+2}{x-2}= \lim\limits_{x\to 2}\dfrac{(x-2)(x-1)}{x-2}= \lim\limits_{x\to 2}(x-1)=2-1=1$.
			\itemch Ta có $\lim\limits_{x\to-1}\dfrac{2x^2+3x+1}{x^2-1}= \lim\limits_{x\to-1}\dfrac{(2x+1)(x+1)}{(x-1)(x+1)}= \lim\limits_{x\to-1}\dfrac{2x+1}{x-1}=\dfrac{-2+1}{-1-1}=\dfrac{1}{2}$.
		\end{itemchoice}
	}
\end{ex}

\begin{ex}%[1D3H2-7]%[Dự án tex tài liệu Dương Hùng-Ngô Tất Thành]
	Cho hàm số $f(x)=\heva{
		&{x-2}  &\text{khi} \, x <-1\\
		&\sqrt{x^2+1} &\text{khi} \, x\ge-1}$. Khi đó.
	\choiceTF
	{Giới hạn$\lim\limits_{x\to-2}f(x)=\sqrt 5$}
	{\True Giới hạn$\lim\limits_{x\to-1^-}f(x)=-3$}
	{\True Giới hạn$\lim\limits_{x\to-1^+}f(x)=\sqrt 2$}
	{Hàm số tồn tại giới hạn khi $ x\to-1$}
	\loigiai{
		\begin{itemchoice}
			\itemch Ta có $\lim\limits_{x\to-2}f(x)=\lim\limits_{x\to-2}(x-2)=-4$
			\itemch Xét dãy số $(x_n)$ bất kì sao cho $x_n<-1$ và $x_n\to-1$, ta có $f(x_n)=x_n-2$.\\
			Khi đó $\lim\limits_{x\to-1^-}f(x)=\lim f(x_n)=-1-2=-3$.
			\itemch Xét dãy số $(x_n)$ bất kì sao cho $x_n>-1$ và $x_n\to-1$, ta có $f(x_n)=\sqrt{x_n^2+1}$.\\
			Khi đó $\lim\limits_{x\to-1^+}f(x)=\lim f(x_n)=\sqrt{(-1)^2+1}=\sqrt 2$.
			\itemch $\lim\limits_{x\to-1^-}f(x)\ne \lim\limits_{x\to-1^+}f(x)$ (hay $-3\ne \sqrt 2$) nên không tồn tại $\lim\limits_{x\to -1}f(x)$.
		\end{itemchoice}
	}
\end{ex}

\begin{ex}%[1D3H2-3]%[Dự án tex tài liệu Dương Hùng-Ngô Tất Thành]
	Xét tính đúng sai trong các mệnh đề sau
	\choiceTF
	{\True $\lim\limits_{x\to 0}(-5x^3-4x+2)=2$}
	{$\lim\limits_{x\to-1}\dfrac{2x-3x^2}{4x+1}=-\dfrac{3}{4}$}
	{$\lim\limits_{x\to-5}\dfrac{x^2+2x-15}{x+5}=+\infty $}
	{\True $\lim\limits_{x\to-4}\dfrac{x^2+3x-4}{x^2+4x}=\dfrac{5}{4}$}
	\loigiai{
		\begin{itemchoice}
			\itemch $\lim\limits_{x\to 0}(-5x^3-4x+2)=-5\cdot{0^3}-4\cdot 0+2=2$.
			\itemch $\lim\limits_{x\to-1}\dfrac{2x-3x^2}{4x+1}=\dfrac{2\cdot (-1)-3\cdot{(-1)^2}}{4\cdot (-1)+1}=\dfrac{5}{3}$.
			\itemch $\lim\limits_{x\to-5}\dfrac{x^2+2x-15}{x+5}=\lim\limits_{x\to-5}\dfrac{(x+5)(x-3)}{x+5}=\lim\limits_{x\to-5}(x-3)=-5-3=-8$.
			\itemch $\lim\limits_{x\to-4}\dfrac{x^2+3x-4}{x^2+4x}= \lim\limits_{x\to-4}\dfrac{(x-1)(x+4)}{x(x+4)}= \lim\limits_{x\to-4}\dfrac{x-1}{x}=\dfrac{-4-1}{-4}=\dfrac{5}{4}$.
		\end{itemchoice}
	}
\end{ex}

\begin{ex}%[1D3H2-7]%[Dự án tex tài liệu Dương Hùng-Ngô Tất Thành]
	Cho hàm số $f(x)=\heva{
		&{1-x^2} &\text{khi} \,x < 2\\
		&\sqrt{x+2} &\text{khi} \,x\ge 2}$
	\choiceTF
	{Giới hạn $\lim\limits_{x\to 3}f(x)=-8$}
	{\True Giới hạn $\lim\limits_{x\to{2^-}}f(x)=-3$}
	{\True Giới hạn $\lim\limits_{x\to{2^+}}f(x)=2$}
	{Giới hạn $\lim\limits_{x\to 2}f(x)=4$}
	\loigiai{
		\begin{itemchoice}
			\itemch Ta có $\lim\limits_{x\to 3}f(x)=\sqrt 5$.
			\itemch Xét dãy số $(x_n)$ bất kì sao cho $x_n< 2$ và $x_n\to 2$, ta có $ f(x_n)=1-x_n^2$.\\
			Khi đó $\lim\limits_{x\to{2^-}}f(x)=\lim f\left(x_n\right)=1-2^2=-3$.
			\itemch Xét dãy số $(x_n)$ bất kì sao cho $x_n> 2$ và $x_n\to 2$, ta có $ f(x_n)=\sqrt{x_n+2}$.\\
			Khi đó $\lim\limits_{x\to{2^+}}f(x)=\lim f\left(x_n\right)=\sqrt{2+2}=2$.
			\itemch Vì $\lim\limits_{x\to{2^-}}f(x)\ne\lim\limits_{x\to{2^+}}f(x)$ (hay $-3\ne 2$) nên không tồn tại $\lim\limits_{x\to 2}f(x)$.
		\end{itemchoice}
	}
\end{ex}
\Closesolutionfile{ans}
\indapan{3}{ans/ans-1-C3B2-DS}

\Opensolutionfile{ans}[ans/ans-1-C3B2-KQ]
\caukq
\begin{ex}%[1D3V2-5]%[Dự án tex tài liệu Dương Hùng-Ngô Tất Thành]
	Tìm giới hạn $\lim\limits_{x\to+\infty}\left(\sqrt{x^2-x}-\sqrt[3]{x^3+1}\right)$.
	\shortans{$-0{,}5$}
	\loigiai{
		\begin{eqnarray*}
			\lim\limits_{x\to+\infty}\left(\sqrt{x^2-x}-\sqrt[3]{x^3+1}\right)&=& \lim\limits_{x\to+\infty}\left(\sqrt{x^2-x}-x\right)+\lim\limits_{x\to+\infty}\left(x-\sqrt[3]{x^3+1}\right)\\
			&=&\lim\limits_{x\to+\infty}\dfrac{x^2-x-x^2}{\sqrt{x^2-x}+x}+\lim\limits_{x\to+\infty}\dfrac{x^3-x^3-1}{x^2+x\sqrt[3]{x^3+1}+\sqrt[3]{\left(x^3+1\right)^2}}\\
			&=&\lim\limits_{x\to+\infty}\dfrac{-x}{\sqrt{x^2-x}+x}+\lim\limits_{x\to+\infty}\dfrac{-1}{x^2+x\sqrt[3]{x^3+1}+\sqrt[3]{\left(x^3+1\right)^2}}
		\end{eqnarray*}
		Xét $A=\lim\limits_{x\to+\infty}\dfrac{-x}{\sqrt{x^2-x}+x}= \lim\limits_{x\to+\infty}\dfrac{-x}{\sqrt{x^2\left(1-\dfrac{1}{x}\right)}+x}\\
		=\lim\limits_{x\to+\infty}\dfrac{-x}{x\sqrt{1-\dfrac{1}{x}}+x}=\lim\limits_{x\to+\infty}\dfrac{-1}{\sqrt 1+1}=-\dfrac{1}{2}$.\\
		Xét $B=\lim\limits_{x\to+\infty}\dfrac{-1}{x^2+x\sqrt[3]{x^3+1}+\sqrt[3]{\left(x^3+1\right)^2}}=0$,\\
		do $\lim\limits_{x\to+\infty}\left(x^2+x\sqrt[3]{x^3+1}+\sqrt[3]{\left(x^3+1\right)^2}\right)=+\infty $\\
		Suy ra $\lim\limits_{x\to+\infty}\left(\sqrt{x^2-x}-\sqrt[3]{x^3+1}\right)=A+B=-\dfrac{1}{2}=-0{,}5$.
	}
\end{ex}

\begin{ex}%%[1D3V2-5]%[Dự án tex tài liệu Dương Hùng-Ngô Tất Thành]
	Tìm giới hạn $\lim\limits_{x\to 0}\dfrac{\sqrt[3]{1+x^2}-1}{x^2}$ (kết quả làm tròn đến hàng phần trăm).
	\shortans{$0{,}33$}
	\loigiai{
		\begin{eqnarray*}
			\lim\limits_{x\to 0}\dfrac{\sqrt[3]{1+x^2}-1}{x^2}&=&\lim\limits_{x\to 0}\dfrac{1+x^2-1^3}{\left(\sqrt[3]{\left(1+x^2\right)^2}+\sqrt[3]{1+x^2}+1\right){x^2}}\\
			&=&\lim\limits_{x\to 0}\dfrac{1}{\sqrt[3]{(1+x^2)^2}+\sqrt[3]{1+x^2}+1}=\dfrac{1}{3} \approx 0{,}33.
		\end{eqnarray*}
	}
\end{ex}

\begin{ex}%[1D3V2-5]%[Dự án tex tài liệu Dương Hùng-Ngô Tất Thành]
	Tìm giới hạn $\lim\limits_{x\to 1}\dfrac{\sqrt[3]{x+7}-\sqrt{x+3}}{x^2-3x+2}$ (kết quả làm tròn đến hàng phần trăm).
	\shortans{$0{,}17$}
	\loigiai{
		$\lim\limits_{x\to 1}\dfrac{\sqrt[3]{x+7}-\sqrt{x+3}}{x^2-3x+2}= \lim\limits_{x\to 1}\dfrac{\sqrt[3]{x+7}-2}{(x-1)(x-2)}+\lim\limits_{x\to 1}\dfrac{2-\sqrt{x+3}}{(x-1)(x-2)}$.\\
		Ta có 
		\begin{eqnarray*}
			\lim\limits_{x\to 1}\dfrac{\sqrt[3]{x+7}-2}{(x-1)(x-2)}&=& \lim\limits_{x\to 1}\dfrac{x+7-2^3}{(x-1)(x-2)\left[\sqrt[3]{(x+7)^2}+2\sqrt[3]{x+7}+4\right]}\\
			&=&\lim\limits_{x\to 1}\dfrac{1}{(x-2)\left(\sqrt[3]{(x+7)^2}+2\sqrt[3]{x+7}+4\right)}= -\dfrac{1}{12}.\\
			\lim\limits_{x\to 1}\dfrac{2-\sqrt{x+3}}{(x-1)(x-2)}&=& \lim\limits_{x\to 1}\dfrac{2^2-(x+3)}{(x-1)(x-2)(2+\sqrt{x+3})}\\
			&=& \lim\limits_{x\to 1}\dfrac{-1}{(x-2)(2+\sqrt{x+3})}=\dfrac{1}{4}.
		\end{eqnarray*}
		Vậy $\lim\limits_{x\to 1}\dfrac{\sqrt[3]{x+7}-\sqrt{x+3}}{x^2-3x+2}= -\dfrac{1}{12}+\dfrac{1}{4}=\dfrac{1}{6} \approx 0{,}17$.
	}
\end{ex}

\begin{ex}%[1D3V2-5]%[Dự án tex tài liệu Dương Hùng-Ngô Tất Thành]
	Tính giới hạn $\lim\limits_{x\to 1}\dfrac{\sqrt{3x+1}-2}{2x^2-3x+1}$
	\shortans{$0{,}75$}
	\loigiai{
		\begin{eqnarray*}
		\lim\limits_{x\to 1}\dfrac{\sqrt{3x+1}-2}{2x^2-3x+1}
		&=&\lim\limits_{x\to 1}\dfrac{3x+1-4}{\left(2x^2-3x+1\right)(\sqrt{3x+1}+2)}\\
		&=&\lim\limits_{x\to 1}\dfrac{3(x-1)}{(x-1)(2x-1)(\sqrt{3x+1}+2)}\\
		&=&\lim\limits_{x\to 1}\dfrac{3}{(2x-1)(\sqrt{3x+1}+2)}=\dfrac{3}{4} = 0{,}75.
	\end{eqnarray*}
	}
\end{ex}

\begin{ex}%[1D3V2-5]%[Dự án tex tài liệu Dương Hùng-Ngô Tất Thành]
	Tính giới hạn $\lim\limits_{x\to-1}\dfrac{3x+\sqrt{4-5x}}{-2x^2+3x+5}$ kết quả làm tròn đến hàng phần trăm.
	\shortans{$0{,}31$}
	\loigiai{
		\begin{eqnarray*}
		\lim\limits_{x\to-1}\dfrac{3x+\sqrt{4-5x}}{-2x^2+3x+5}
		&=&\lim\limits_{x\to-1}\dfrac{9x^2-4+5x}{\left(-2x^2+3x+5\right)(3x-\sqrt{4-5x})}\\
		&=&\lim\limits_{x\to-1}\dfrac{(x+1)(9x-4)}{(x+1)(-2x+5)(3x-\sqrt{4-5x})}\\
		&=&\lim\limits_{x\to-1}\dfrac{9x-4}{(-2x+5)(3x-\sqrt{4-5x})}=\dfrac{13}{42} \approx 0{,}31.
		\end{eqnarray*}
	}
\end{ex}

\begin{ex}%[1D3V2-5]%[Dự án tex tài liệu Dương Hùng-Ngô Tất Thành]
	Tính giới hạn $\lim\limits_{x\to 1}\dfrac{\sqrt{x+3}+\sqrt{2x+7}-5}{2x-2}$ kết quả làm tròn đến hàng phần trăm.
	\shortans{$0{,}29$}
	\loigiai{
	\begin{eqnarray*}
		\lim\limits_{x\to 1}\dfrac{\sqrt{x+3}+\sqrt{2x+7}-5}{2x-2}
		&=&\lim\limits_{x\to 1}\dfrac{\sqrt{x+3}-2+\sqrt{2x+7}-3}{2x-2}\\
		&=&\lim\limits_{x\to 1}\left[\dfrac{x+3-4}{(2x-2)(\sqrt{x+3}+2)}+\dfrac{2x+7-9}{(2x-2)(\sqrt{2x+7}+3)}\right]\\
		&=&\lim\limits_{x\to 1}\left[\dfrac{x-1}{2(x-1)(\sqrt{x+3}+2)}+\dfrac{2x-2}{(2x-2)(\sqrt{2x+7}+3)}\right]\\
		&=&\lim\limits_{x\to 1}\left[\dfrac{1}{2(\sqrt{x+3}+2)}+\dfrac{1}{\sqrt{2x+7}+3}\right]
		=\dfrac{7}{24} \approx 0{,}29.
	\end{eqnarray*}
	}
\end{ex}
\Closesolutionfile{ans}
\indapan{6}{ans/ans-1-C3B2-KQ}

\subsection{Đề 2}
\caulc
\Opensolutionfile{ans}[ans/ans-1-C3B2]
\begin{ex}%[1D3H2-7]%[Dự án tex tài liệu Dương Hùng-Ngô Tất Thành]
		Kết quả của giới hạn  $\lim\limits_{x\to{\left(-2\right)^+}}\dfrac{|3x+6|}{x+2}$ là
	\choice
		{$-\infty$}
		{\True $3$}
		{$+\infty$}
		{Không xác định}
	\loigiai{
		Ta có $|x+2|=x+2$ với mọi $x >-2$,\\
		do đó $\lim\limits_{x\to{(-2)^+}}\dfrac{|3x+6|}{x+2} =\lim\limits_{x\to{(-2)^+}}\dfrac{3|x+2|}{x+2} =\lim\limits_{x\to{(-2)^+}}\dfrac{3(x+2)}{x+2} =\lim\limits_{x\to{(-2)^+}}3=3$.
	}
\end{ex}

\begin{ex}%[1D3H2-7]%[Dự án tex tài liệu Dương Hùng-Ngô Tất Thành]
		Kết quả của giới hạn  $\lim\limits_{x\to{2^-}}\dfrac{|2-x|}{2x^2-5x+2}$ là
	\choice
		{$-\infty$}
		{$+\infty$}
		{\True $-\dfrac{1}{3}$}
		{$\dfrac{1}{3}$}
	\loigiai{
		Ta có $\lim\limits_{x\to{2^-}}\dfrac{|2-x|}{2x^2-5x+2} =\lim\limits_{x\to{2^-}}\dfrac{2-x}{(2-x)(1-2x)} =\lim\limits_{x\to{2^-}}\dfrac{1}{1-2x}=-\dfrac{1}{3}$.
	}
\end{ex}

\begin{ex}%[1D3H2-7]%[Dự án tex tài liệu Dương Hùng-Ngô Tất Thành]
		Kết quả của giới hạn  $\lim\limits_{x\to-3^+}\dfrac{x^2+13x+30}{\sqrt{(x+3)(x^2+5)}}$ là
	\choice
		{$-2$}
		{$2$}
		{\True $0$}
		{$\dfrac{2}{\sqrt{15}}$}
	\loigiai{
		Ta có $x+3 > 0$ với mọi $x >-3$, nên
	\begin{eqnarray*}
			\lim\limits_{x\to-3^+}\dfrac{x^2+13x+30}{\sqrt{(x+3)(x^2+5)}} &=&\lim\limits_{x\to-3^+}\dfrac{(x+3)(x+10)}{\sqrt{(x+3)(x^2+5)}}\\ &=&\lim\limits_{x\to-3^+}\dfrac{\sqrt{x+3}(x+10)}{\sqrt{x^2+5}}\\ &=&\dfrac{\sqrt{-3+3}(-3+7)}{\sqrt{(-3)^2+5}}=0.
	\end{eqnarray*}
	}
\end{ex}

\begin{ex}%[1D3H2-7]%[Dự án tex tài liệu Dương Hùng-Ngô Tất Thành]
		Cho hàm số $f(x)=\heva{
			&\dfrac{2x}{\sqrt{1-x}} & \text{khi} \,  x < 1\\
			&\sqrt{3x^2+1} & \text{khi} \,  x\ge 1}$. Khi đó  $\lim\limits_{x\to{1^+}}f(x)$ là
	\choice
		{$+\infty$}
		{\True $2$}
		{$4$}
		{$-\infty$}
	\loigiai{
		$\lim\limits_{x\to{1^+}}f(x)=\lim\limits_{x\to{1^+}}\sqrt{3x^2+1} =\sqrt{3 \cdot 1^2+1}=2$.
	}
\end{ex}

\begin{ex}%[1D3H2-7]%[Dự án tex tài liệu Dương Hùng-Ngô Tất Thành]
		Cho hàm số $f(x)=\heva{
			&\dfrac{x^2+1}{1-x} & \text{khi} \, x < 1\\
			&\sqrt{2x-2}  & \text{khi} \, x\ge 1}$. Khi đó  $\lim\limits_{x\to{1^-}}f(x)$ là
	\choice
		{\True $+\infty $}
		{$-1$}
		{$0$}
		{$1$}
	\loigiai{
		$\lim\limits_{x\to{1^-}}f(x)=\lim\limits_{x\to{1^-}}\dfrac{x^2+1}{1-x} =+\infty$\\
		vì $\heva{
			&\lim\limits_{x\to{1^-}}(x^2+1)=2\\
			&\lim\limits_{x\to{1^-}}(1-x)=0\\
			&1-x > 0, \,(\forall x < 1).}$
		}
\end{ex}

\begin{ex}%[1D3H2-7]%[Dự án tex tài liệu Dương Hùng-Ngô Tất Thành]
		Cho hàm số $f(x)=\heva{
				&x^2-3 & \text{khi} \, x\ge 2\\
				&{x-1} & \text{khi} \, x < 2}$. Khi đó $\lim\limits_{x\to 2}f(x)$ là
	\choice
		{$-1$}
		{$0$}
		{\True $1$}
		{Không tồn tại}
	\loigiai{
		Ta có $\heva{
			&\lim\limits_{x\to{2^+}}f(x)=\lim\limits_{x\to{2^+}}\left(x^2-3\right)=1\\
			&\lim\limits_{x\to{2^-}}f(x)=\lim\limits_{x\to{2^-}}\left(x-1\right)=1} \Rightarrow \lim\limits_{x\to{2^+}}f(x)=\lim\limits_{x\to{2^-}}f(x)=1 \Rightarrow \lim\limits_{x\to 2}f(x)=1$.
		}
\end{ex}

\begin{ex}%[1D3H2-7]%[Dự án tex tài liệu Dương Hùng-Ngô Tất Thành]
		Cho hàm số $f(x)=\heva{
			&\sqrt{x-2}+3 & \text{khi} \,x\ge 2\\
			&ax-1 & \text{khi} \, x < 2}$. Tìm $a$ để tồn tại $\lim\limits_{x\to 2}f(x)$.
	\choice
		{$a=1$}
		{\True $a=2$}
		{$a=3$}
		{$a=4$}
	\loigiai{
		Ta có $\heva{
			&\lim\limits_{x\to{2^-}}f(x)=\lim\limits_{x\to{2^-}}\left(ax-1\right)=2a-1\\
			&\lim\limits_{x\to{2^+}}f(x)=\lim\limits_{x\to{2^+}}\left(\sqrt{x-2}+3\right)=3.}$\\
		Khi đó $\lim\limits_{x\to 2}f(x)$ tồn tại $\Leftrightarrow\lim\limits_{x\to{2^-}}f(x)=\lim\limits_{x\to{2^+}}f(x) \Leftrightarrow 2a-1=3\Leftrightarrow a=2$.
	}
\end{ex}

\begin{ex}%[1D3H2-7]%[Dự án tex tài liệu Dương Hùng-Ngô Tất Thành]
		Cho hàm số $f(x)=\heva{
			&x^2-2x+3 &\text{khi} \, x > 3\\
			&1 &\text{khi} \, x=3\\
			&3-2x^2 &\text{khi} khi \, x < 3}$. Khẳng định nào dưới đây \textbf{sai}?
	\choice
		{$\lim\limits_{x\to{3^+}}f(x)=6$}
		{$f(3)=1$}
		{\True $\lim\limits_{x\to{3}}f(x)=6$}
		{$\lim\limits_{x\to{3^-}}f(x)=-15$}
	\loigiai{
		Ta có $\heva{
			&\lim\limits_{x\to{3^+}}f(x)=\lim\limits_{x\to{3^+}}(x^2-2x+3)=6\\
			&\lim\limits_{x\to{3^-}}f(x)=\lim\limits_{x\to{3^-}}(3-2x^2)=-15} \Rightarrow \lim\limits_{x\to{3^+}}f(x) \ne \lim\limits_{x\to{3^-}}f(x)$\\
		Suy ra không tồn tại giới hạn khi $x\to 3$.\\
	}
\end{ex}

\begin{ex}%[1D3H2-4]%[Dự án tex tài liệu Dương Hùng-Ngô Tất Thành]
		Giá trị của giới hạn  $\lim\limits_{x\to+\infty}\left(\sqrt{1+2x^2}-x\right)$ là
	\choice
		{$0$}
		{\True $+\infty $}
		{$\sqrt 2-1$}
		{$-\infty $}
	\loigiai{
		Ta có  $\lim\limits_{x\to+\infty}\left(\sqrt{1+2x^2}-x\right) =\lim\limits_{x\to+\infty}x\left(\sqrt{\dfrac{1}{x^2}+2}-1\right) =+\infty$\\
		Vì $\lim\limits_{x\to+\infty}x=+\infty$, $\lim\limits_{x\to+\infty}\left(\sqrt{\dfrac{1}{x^2}+2}-1\right)=\sqrt 2-1 > 0$.
	}
\end{ex}

\begin{ex}%[1D3H2-5]%[Dự án tex tài liệu Dương Hùng-Ngô Tất Thành]
	Giá trị của giới hạn  $\lim\limits_{x\to+\infty}\left(\sqrt{x^2+1}-x\right)$ là
	\choice
		{\True $0$}
		{$+\infty $}
		{$\dfrac{1}{2}$}
		{$-\infty $}
	\loigiai{
		$\lim\limits_{x\to+\infty}\left(\sqrt{x^2+1}-x\right)= \lim\limits_{x\to+\infty}\dfrac{1}{\sqrt{x^2+1}+x}= \lim\limits_{x\to+\infty}\dfrac{\dfrac{1}{x}}{\sqrt{1+\dfrac{1}{x^2}}+1}= \dfrac{0}{2}=0$.
	}
\end{ex}

\begin{ex}%[1D3H2-5]%[Dự án tex tài liệu Dương Hùng-Ngô Tất Thành]
		Giá trị của giới hạn  $\lim\limits_{x\to+\infty}(\sqrt{x^2+3x}-\sqrt{x^2+4x})$ là
	\choice
		{$\dfrac{7}{2}$}
		{\True $-\dfrac{1}{2}$}
		{$+\infty $}
		{$-\infty $}
	\loigiai{
		\begin{eqnarray*}
		\lim\limits_{x\to+\infty}(\sqrt{x^2+3x}-\sqrt{x^2+4x}) &=&\lim\limits_{x\to+\infty}\dfrac{-x}{\sqrt{x^2+3x}+\sqrt{x^2+4x}}\\
		&=&\lim\limits_{x\to+\infty}\dfrac{-1}{\sqrt{1+\dfrac{3}{x}}+\sqrt{1+\dfrac{4}{x}}}=-\dfrac{1}{2}.
		\end{eqnarray*}
	}
\end{ex}

\begin{ex}%[1D3V2-5]%[Dự án tex tài liệu Dương Hùng-Ngô Tất Thành]
		Giá trị của giới hạn  $\lim\limits_{x\to+\infty}\left(\sqrt{x^2+x}-\sqrt[3]{x^3-x^2}\right)$ là
	\choice
		{\True $\dfrac{5}{6}$}
		{$+\infty $}
		{$-1$}
		{$-\infty $}
	\loigiai{
		\begin{eqnarray*}
		\lim\limits_{x\to+\infty}\left(\sqrt{x^2+x}-\sqrt[3]{x^3-x^2}\right)&=& \lim\limits_{x\to+\infty}\left(\sqrt{x^2+x}-x+x-\sqrt[3]{x^3-x^2}\right)\\
		&=&\lim\limits_{x\to+\infty}\left(\dfrac{x}{\sqrt{x^2+1}+x}+\dfrac{x^2}{x^2+x\sqrt[3]{x^3-1}+\sqrt[3]{\left(x^3-1\right)^2}}\right)\\ &=&\dfrac{1}{2}+\dfrac{1}{3}=\dfrac{5}{6}.
		\end{eqnarray*}
	}
\end{ex}
\Closesolutionfile{ans}
\indapan{6}{ans/ans-1-C3B2}

\cauds
\Opensolutionfile{ans}[ans/ans-1-C3B2-DS]

\begin{ex}%[1D3H2-3]%[Dự án tex tài liệu Dương Hùng-Ngô Tất Thành]
		Tìm được các giới hạn sau.
	\choiceTF
		{\True $\lim\limits_{x\to+\infty}\left(x^2+3\right)=+\infty$}
		{$\lim\limits_{x\to-\infty}\left(\sqrt{x^2+x}-x\right)=-\infty $}
		{\True $\lim\limits_{x\to-\infty}\dfrac{1}{x+2}=0$}
		{$\lim\limits_{x\to+\infty}\sqrt{\dfrac{2x}{x+3}}=2$}
	\loigiai{
	\begin{itemchoice}
		\itemch $\lim\limits_{x\to+\infty}\left(x^2+3\right)= \lim\limits_{x\to+\infty}{x^2}\left(1+\dfrac{3}{x^2}\right)=+\infty$, \\
		do $\lim\limits_{x\to+\infty}{x^2}=+\infty$ và $\lim\limits_{x\to+\infty}\left(1+\dfrac{3}{x^2}\right)=1$.
		\itemch $\lim\limits_{x\to-\infty}\left(\sqrt{x^2+x}-x\right)= \lim\limits_{x\to-\infty}\left(-x\sqrt{1+\dfrac{1}{x}}-x\right)\\
		=\lim\limits_{x\to-\infty}x\left(-\sqrt{1+\dfrac{1}{x}}-1\right)=+\infty$, \\
		do $\lim\limits_{x\to-\infty}x=-\infty$ và $\lim\limits_{x\to-\infty}\left(-\sqrt{1+\dfrac{1}{x}}-1\right)=-2$.
		\itemch $\lim\limits_{x\to-\infty}\dfrac{1}{x+2}= \lim\limits_{x\to-\infty}\dfrac{x\cdot\dfrac{1}{x}}{x\left(1+\dfrac{2}{x}\right)}= \lim\limits_{x\to-\infty}\dfrac{\dfrac{1}{x}}{1+\dfrac{2}{x}}=0$.
		\itemch $\lim\limits_{x\to+\infty}\sqrt{\dfrac{2x}{x+3}}= \lim\limits_{x\to+\infty}\sqrt{\dfrac{2x}{x\left(1+\dfrac{3}{x}\right)}}= \lim\limits_{x\to+\infty}\sqrt{\dfrac{2}{1+\dfrac{3}{x}}}=\sqrt 2 $.
	\end{itemchoice}
	}
\end{ex}
	
\begin{ex}%[1D3H2-4]%[Dự án tex tài liệu Dương Hùng-Ngô Tất Thành]
		Tìm được các giới hạn sau.
	\choiceTF
		{\True $\lim\limits_{x\to{2^+}}(\sqrt{x+2}-1)=1$}
		{\True $\lim\limits_{x\to{1^+}}\dfrac{4x-3}{x-1}=+\infty $}
		{\True $\lim\limits_{x\to{2^-}}\left(\dfrac{1}{x-2}-\dfrac{1}{x^2-4}\right)=-\infty $}
		{$\lim\limits_{x\to-1^-}\dfrac{|x+1|}{x^2-1}=-\infty $}
	\loigiai{
	\begin{itemchoice}
		\itemch $\lim\limits_{x\to{2^+}}(\sqrt{x+2}-1)=\sqrt{2+2}-1=1$.
		\itemch $\lim\limits_{x\to{1^+}}\dfrac{4x-3}{x-1}= \lim\limits_{x\to{1^+}}\left[(4x-3)\cdot\dfrac{1}{x-1}\right]=+\infty$ \\
		vì $\lim\limits_{x\to{1^+}}(4x-3)=1$, $\lim\limits_{x\to{1^+}}\dfrac{1}{x-1}=+\infty $.
		\itemch \begin{eqnarray*}
			\lim\limits_{x\to{2^-}}\left(\dfrac{1}{x-2}-\dfrac{1}{x^2-4}\right)&=& \lim\limits_{x\to{2^-}}\dfrac{x+2-1}{(x-2)(x+2)}\\
			&=& \lim\limits_{x\to{2^-}}\dfrac{x+1}{(x-2)(x+2)}\\ &=&\lim\limits_{x\to{2^-}}\left[\dfrac{x+1}{x+2}\cdot\dfrac{1}{(x-2)}\right]= -\infty,
		\end{eqnarray*}
		do $\heva{
						&\lim\limits_{x\to{2^-}}\dfrac{x+1}{x+2}=\dfrac{3}{4}\\
						&\lim\limits_{x\to{2^-}}\dfrac{1}{x-2}=-\infty}$
		\itemch $\lim\limits_{x\to-1^-}\dfrac{|x+1|}{x^2-1}=\lim\limits_{x\to-1^-}\dfrac{-x-1}{(x-1)(x+1)}=\lim\limits_{x\to-1^-}\dfrac{-1}{x-1}=\dfrac{1}{2}$.
	\end{itemchoice}
	}
\end{ex}

\begin{ex}%[1D3H2-5]%[Dự án tex tài liệu Dương Hùng-Ngô Tất Thành]
		Tìm được các giới hạn sau.
	\choiceTF
		{\True $\lim\limits_{x\to 0}\dfrac{\sqrt{4+x}-2}{4x}=\dfrac{1}{16}$}
		{\True $\lim\limits_{x\to 2}\dfrac{4-x^2}{\sqrt{x+7}-3}=-24$}
		{\True $\lim\limits_{x\to 2}\dfrac{\sqrt{2x+5}-3}{\sqrt{x+2}-2}=\dfrac{4}{3}$}
		{$\lim\limits_{x\to 1}\dfrac{\sqrt[3]{x+7}-2}{x-1}=\dfrac{1}{3}$}
	\loigiai{
		\begin{itemchoice}
			\itemch \begin{eqnarray*}
			\lim\limits_{x\to 0}\dfrac{\sqrt{4+x}-2}{4x}&=& \lim\limits_{x\to 0}\dfrac{(\sqrt{4+x}-2)(\sqrt{4+x}+2)}{4x(\sqrt{4+x}+2)}=\lim\limits_{x\to 0}\dfrac{4+x-4}{4x(\sqrt{4+x}+2)}\\
			&=&\lim\limits_{x\to 0}\dfrac{1}{4(\sqrt{4+x}+2)}= \dfrac{1}{4(\sqrt 4+2)}=\dfrac{1}{16}.
			\end{eqnarray*}
			\itemch \begin{eqnarray*}
			\lim\limits_{x\to 2}\dfrac{4-x^2}{\sqrt{x+7}-3}&=& \lim\limits_{x\to 2}\dfrac{(2-x)(2+x)(\sqrt{x+7}+3)}{(\sqrt{x+7}-3)(\sqrt{x+7}+3)}\\
			&=& \lim\limits_{x\to 2}\dfrac{(2-x)(2+x)(\sqrt{x+7}+3)}{x+7-9}\\
			&=&\lim\limits_{x\to 2}[-(2+x)(\sqrt{x+7}+3)]=-4 \cdot 6=-24.
			\end{eqnarray*}
			\itemch \begin{eqnarray*}
			\lim\limits_{x\to 2}\dfrac{\sqrt{2x+5}-3}{\sqrt{x+2}-2}&=& \lim\limits_{x\to 2}\dfrac{(\sqrt{2x+5}-3)(\sqrt{2x+5}+3)(\sqrt{x+2}+2)}{(\sqrt{x+2}-2)(\sqrt{x+2}+2)(\sqrt{2x+5}+3)}\\
			&=& \lim\limits_{x\to 2}\dfrac{(2x+5-9)(\sqrt{x+2}+2)}{(x+2-4)(\sqrt{2x+5}+3)}\\	&=&\lim\limits_{x\to 2}\dfrac{2(\sqrt{x+2}+2)}{\sqrt{2x+5}+3}=\dfrac{4}{3}.
			\end{eqnarray*}
			\itemch \begin{eqnarray*}
			\lim\limits_{x\to 1}\dfrac{\sqrt[3]{x+7}-2}{x-1}&=& \lim\limits_{x\to 1}\dfrac{(\sqrt[3]{x+7}-2)\left(\sqrt[3]{(x+7)^2}+2\sqrt[3]{x+7}+4\right)}{(x-1)\left(\sqrt[3]{(x+7)^2}+2\sqrt[3]{x+7}+4\right)}\\
			&=& \lim\limits_{x\to 1}\dfrac{x+7-2^3}{(x-1)\left(\sqrt[3]{(x+7)^2}+2\sqrt[3]{x+7}+4\right)}\\
			&=&\lim\limits_{x\to 1}\dfrac{1}{\sqrt[3]{(x+7)^2}+2\sqrt[3]{x+7}+4}=\dfrac{1}{12}.
			\end{eqnarray*}
		\end{itemchoice}
	}
\end{ex}

\begin{ex}%[1D3V2-4]%[Dự án tex tài liệu Dương Hùng-Ngô Tất Thành]
		Tìm được các giới hạn sau.
	\choiceTF
		{\True $\lim\limits_{x\to-\infty}\left(x^2-10x\right)=+\infty $}
		{\True $\lim\limits_{x\to+\infty}\dfrac{3x^2-4x+1}{2x^2+x+1}=\dfrac{3}{2}$}
		{$\lim\limits_{x\to-\infty}\dfrac{\sqrt{x^2+x+1}-3x}{2-3x}=\dfrac{5}{4}$}
		{\True $\lim\limits_{x\to-\infty}\dfrac{\sqrt[3]{8x^3+3x^2+1}-x}{\sqrt{4x^2-x+2}+3x}=1$}
	\loigiai{
	\begin{itemchoice}
		\itemch $\lim\limits_{x\to-\infty}\left(x^2-10x\right)= \lim\limits_{x\to-\infty}{x^2}\left(1-\dfrac{10}{x}\right)=+\infty $.
		\itemch $\lim\limits_{x\to+\infty}\dfrac{3x^2-4x+1}{2x^2+x+1}= \lim\limits_{x\to+\infty}\dfrac{x^2\left(3-\dfrac{4}{x}+\dfrac{1}{x^2}\right)}{x^2\left(2+\dfrac{1}{x}+\dfrac{1}{x^2}\right)}= \lim\limits_{x\to+\infty}\dfrac{3-\dfrac{4}{x}+\dfrac{1}{x^2}}{2+\dfrac{1}{x}+\dfrac{1}{x^2}}=\dfrac{3}{2}$
		\itemch \begin{eqnarray*}
		\lim\limits_{x\to-\infty}\dfrac{\sqrt{x^2+x+1}-3x}{2-3x}&=& \lim\limits_{x\to-\infty}\dfrac{\sqrt{x^2\left(1+\dfrac{1}{x}+\dfrac{1}{x^2}\right)}-3x}{x\left(\dfrac{2}{x}-3\right)}\\
		&=& \lim\limits_{x\to-\infty}\dfrac{-x\sqrt{1+\dfrac{1}{x}+\dfrac{1}{x^2}}-3x}{x\left(\dfrac{2}{x}-3\right)}\\
		&=&\lim\limits_{x\to-\infty}\dfrac{-\sqrt{1+\dfrac{1}{x}+\dfrac{1}{x^2}}-3}{\dfrac{2}{x}-3}=\dfrac{-\sqrt 1-3}{-3}=\dfrac{4}{3}.
		\end{eqnarray*}
		\itemch \begin{eqnarray*}
		\lim\limits_{x\to-\infty}\dfrac{\sqrt[3]{8x^3+3x^2+1}-x}{\sqrt{4x^2-x+2}+3x}&=&\lim\limits_{x\to-\infty}\dfrac{\sqrt[3]{x^3\left(8+\dfrac{3}{x}+\dfrac{1}{x^3}\right)}-x}{\sqrt{x^2\left(4-\dfrac{1}{x}+\dfrac{2}{x^2}\right)}+3x}\\
		&=&\lim\limits_{x\to-\infty}\dfrac{x\sqrt[3]{8+\dfrac{3}{x}+\dfrac{1}{x^3}}-x}{x\sqrt{4-\dfrac{1}{x}+\dfrac{2}{x^2}}+3x}\\
		&=&\lim\limits_{x\to-\infty}\dfrac{\sqrt[3]{8+\dfrac{3}{x}+\dfrac{1}{x^3}}-1}{\sqrt{4-\dfrac{1}{x}+\dfrac{2}{x^2}}+3}= \dfrac{\sqrt[3]{8}-1}{-\sqrt 4+3}=1.
		\end{eqnarray*}
	\end{itemchoice}
	}
\end{ex}
\Closesolutionfile{ans}
\indapan{3}{ans/ans-1-C3B2-DS}

\caukq
\Opensolutionfile{ans}[ans/ans-1-C3B2-KQ]

\begin{ex}%[1D3H2-5]%[Dự án tex tài liệu Dương Hùng-Ngô Tất Thành]
		Tính giới hạn $\lim\limits_{x\to 2}\dfrac{\sqrt{5x-1}-\sqrt{9x-2}+1}{x-2}$ kết quả làm tròn đến hàng phần mười.
	\shortans{$-0{,}3$}
	\loigiai{
		\begin{eqnarray*}
		\lim\limits_{x\to 2}\dfrac{\sqrt{5x-1}-\sqrt{9x-2}+1}{x-2}
		&=&\lim\limits_{x\to 2}\dfrac{\sqrt{5x-1}-3+4-\sqrt{9x-2}}{x-2}\\ &=&\lim\limits_{x\to 2}\left(\dfrac{\sqrt{5x-1}-3}{x-2}+\dfrac{4-\sqrt{9x-2}}{x-2}\right)\\ 
		&=&\lim\limits_{x\to 2}\left[\dfrac{5x-1-9}{(x-2)(\sqrt{5x-1}+3)}+\dfrac{16-9x+2}{(x-2)(4+\sqrt{9x-2})}\right]\\
		&=&\lim\limits_{x\to 2}\left[\dfrac{5(x-2)}{(x-2)(\sqrt{5x-1}+3)}+\dfrac{9(2-x)}{(x-2)(4+\sqrt{9x-2})}\right]\\
		&=&\lim\limits_{x\to 2}\left(\dfrac{5}{\sqrt{5x-1}+3}-\dfrac{9}{4+\sqrt{9x-2}}\right)=-\dfrac{7}{24} \approx -0{,}3.
		\end{eqnarray*}
	}
\end{ex}
	
\begin{ex}%[1D3V2-3]%[Dự án tex tài liệu Dương Hùng-Ngô Tất Thành]
		Tính giới hạn  $\lim\limits_{x\to+\infty}\dfrac{(3-2x)^5\cdot\left(3x^2-2x\right)}{\left(2x^2-1\right)^2(7-2x)^3}$
	\shortans{$3$}
	\loigiai{
		\begin{eqnarray*}
		\lim\limits_{x\to+\infty}\dfrac{(3-2x)^5\cdot\left(3x^2-2x\right)}{\left(2x^2-1\right)^2(7-2x)^3}&=& \lim\limits_{x\to+\infty}\dfrac{x^5\left(\dfrac{3}{x}-2\right)^5\cdot{x^2}\left(3-\dfrac{2}{x}\right)}{x^4\left(2-\dfrac{1}{x^2}\right)^2\cdot{x^3}{\left(\dfrac{7}{x}-2\right)^3}}\\
		&=&\lim\limits_{x\to+\infty}\dfrac{\left(\dfrac{3}{x}-2\right)^5\cdot\left(3-\dfrac{2}{x}\right)}{\left(2-\dfrac{1}{x^2}\right)^2\cdot{\left(\dfrac{7}{x}-2\right)^3}}=3.
		\end{eqnarray*}
	}
\end{ex}

\begin{ex}%%[1D3V2-5]%[Dự án tex tài liệu Dương Hùng-Ngô Tất Thành]
		Tìm giới hạn hàm số $\lim\limits_{x\to 2}\dfrac{2-\sqrt{x+2}}{x^2-3x+2}$ kết quả dạng phân số $-\dfrac{a}{b}$. Hỏi tổng $2a+3b$ bằng
	\shortans{$14$}
	\loigiai{
		Ta có $\lim\limits_{x\to 2}\dfrac{2-\sqrt{x+2}}{x^2-3x+2}=\dfrac{0}{0}$, để khử dạng vô định ta nhân tử và mẫu cho $(2+\sqrt{x+2})$, ta được\\
		 \begin{eqnarray*}
		 \lim\limits_{x\to 2}\dfrac{2-\sqrt{x+2}}{x^2-3x+2}&=&\lim\limits_{x\to 2}\dfrac{(2-\sqrt{x+2})(2+\sqrt{x+2})}{(x-1)(x-2)(2+\sqrt{x+2})}\\
		 &=&\lim\limits_{x\to 2}\dfrac{-1}{(x-1)(2+\sqrt{x+2})}=-\dfrac{1}{4}.\\
	\end{eqnarray*}
	Vậy tổng $2a+3b=2 \cdot 1 +3 \cdot 4=14$.
	}
\end{ex}

\begin{ex}%%[1D3V2-5]%[Dự án tex tài liệu Dương Hùng-Ngô Tất Thành]
		Tính giới hạn của hàm số sau $\lim\limits_{x\to-2}\dfrac{\sqrt[3]{2x+12}+x}{x^2+2x}$ kết quả dạng phân số $-\dfrac{a}{b}$. Hỏi tổng $a+2b$ bằng
	\shortans{$17$}
	\loigiai{
		\begin{eqnarray*}
		\lim\limits_{x\to-2}\dfrac{\sqrt[3]{2x+12}+x}{x^2+2x}&=&\lim\limits_{x\to-2}\dfrac{(\sqrt[3]{2x+12}+x)\left(\sqrt[3]{(2x+12)^2}-x\sqrt[3]{2x+12}+x^2\right)}{x(x+2)\left(\sqrt[3]{(2x+12)^2}-x\sqrt[3]{2x+12}+x^2\right)}\\
		&=&\lim\limits_{x\to-2}\dfrac{(x+2)\left(x^2-2x+12\right)}{x(x+2)\left(\sqrt[3]{(2x+12)^2}-x\sqrt[3]{2x+12}+x^2\right)}\\
		&=&\lim\limits_{x\to-2}\dfrac{x^2-2x+12}{\left(\sqrt[3]{(2x+12)^2}-x\sqrt[3]{2x+12}+x^2\right)x}=-\dfrac{5}{6}.
		\end{eqnarray*}
		Vậy tổng $a+2b=5+2\cdot 6=17$.
	}
\end{ex}

\begin{ex}%[1D3V2-8]%[Dự án tex tài liệu Dương Hùng-Ngô Tất Thành]
		Một cái hồ chứa $600$ lít nước ngọt. Người ta bơm nước biển có nồng độ muối $30$ gam/lít vào hồ với tốc độ $15$ lít/phút. Nồng độ muối trong hồ dần về bao nhiêu gam/lít khi $t$ dần về dương vô cùng?
	\shortans{$30$}
	\loigiai{
		Sau $t$ phút bơm nước vào hồ thì lượng nước là $600+15t$ (lít) và lượng muối có được là $30 \cdot 15t$ (gam).\\
		Nồng độ muối của nước là $C(t)=\dfrac{30 \cdot 15t}{600+15t}=\dfrac{30t}{40+t}$ (gam/lít).\\
		Khi $t$ dần về dương vô cùng, ta có\\
		$\lim\limits_{t\to+\infty}C(t)=\lim\limits_{t\to+\infty}\dfrac{30t}{40+t}= \lim\limits_{t\to+\infty}\dfrac{30t}{t\left(\dfrac{40}{t}+1\right)}= \lim\limits_{t\to+\infty}\dfrac{30}{\dfrac{40}{t}+1}=30$ (gam/lít).
	}
\end{ex}

\begin{ex}%[1D3C2-8]%[Dự án tex tài liệu Dương Hùng-Ngô Tất Thành]
	\immini
		{Trong hệ trục toạ độ $Oxy$, lấy điểm $A$ thuộc tia $Ox$ và điểm $B(0;2)$ thuộc tia $Oy$. Giả sử hoành độ điểm $A$ là $a > 0$. Khi điểm $A$ dịch chuyển ra vô cực theo chiều dương trục $Ox$ thì độ dài $AH$ dần về bao nhiêu?}
		{
		\begin{tikzpicture}[scale=1,line join=round,line cap=round,font=\footnotesize,>=stealth]
			
			\path 	(0,0) coordinate (O);	
			\draw[-stealth] (0,0)--(4,0)node[above]{$x$};
			\draw[-stealth] (0,0)--(0,3)node[left]{$y$};
			\node[left] at (0,2) {$2$};
			\node[below] at (3,0) {$a$};
			\node[right] at (0,2) {$B$};
			\node[above] at (3,0) {$A$};
			\node[above] at (1,1.5) {$H$};
			\draw (0,2)--(3,0) (0,0)--(0.9,1.4);
		\end{tikzpicture}
		}
	\shortans{$2$}
	\loigiai{
		Tam giác $OAB$ vuông tại $O$ có đường cao $OH$ nên:\\
		$\dfrac{1}{O{H^2}}=\dfrac{1}{O{A^2}}+\dfrac{1}{O{B^2}}\Rightarrow OH= \dfrac{OA\cdot OB}{\sqrt{O{A^2}+O{B^2}}}= \dfrac{2a}{\sqrt{4+a^2}}$.\\
		Đặt $h(a)=OH=\dfrac{2a}{\sqrt{4+a^2}}$.\\
		Khi điểm $A$ dịch chuyển ra vô cực theo tia $Ox$ thì $a\to+\infty $.\\
		Ta có  $\lim\limits_{a\to+\infty}h(a)= \lim\limits_{a\to+\infty}\dfrac{2a}{\sqrt{4+a^2}}= \lim\limits_{a\to+\infty}\dfrac{2a}{a\sqrt{\dfrac{4}{a^2}+1}}= \lim\limits_{a\to+\infty}\dfrac{2}{\sqrt{\dfrac{4}{a^2}+1}}= \dfrac{2}{\sqrt 1}=2$.\\
		Vậy khi điểm $A$ dần về vô cực thì độ dài $OH$ dần về $2$.
	}
\end{ex}
\Closesolutionfile{ans}
\indapan{6}{ans/ans-1-C3B2-KQ}